% discussion.tex -- Discussion
\section{Discussion}

In this study, we explored how a lasting change in the way a minority
of traders formed their intentions could spread selling through a market,
and what then determined how far the price fell. A Sell majority formed
once enough traders received the new response, even though every other trader kept its
original rule (Fig.~\ref{fig:rule-quench}), whereas the size of the
resulting price decline also depended on trader resources and the
dealer's reserve curve. The changed
response acted whenever its neighbor configuration recurred, so the
selected traders did not need to remain in Sell for the change to keep
influencing the market. The tested rule switch produced majorities that a one-time
instruction to sell did not (Supplementary
Table~\ref{tab:revision-pulse}). Testing alternative single-response changes showed that majority
formation depended on which response changed, limiting the conclusions
to the tested rules
(Supplementary Table~\ref{tab:revision-rules}). Because each switched
market shared its history with a no-switch control, the difference
between the two measured the effect of the switch at that starting
market, and repeating the pair across starting markets measured how
often a majority formed within the observation
period~\cite{glasserman1992guidelines,law2015simulation}. These
conclusions concern assigned rule changes and the spread of selling intentions:
traders did not learn the new rule, adopt it for its profitability, or
pass the rule itself to other traders.

The spread of selling did not require any decision to respond to the
price. With the technical traders kept in Hold, trades still completed and the
price still moved, but no decision responded to it, and Sell majorities
still formed (Fig.~\ref{fig:reservoir}). In the displayed markets with active technical traders, a
majority formed as groups of sellers grew and merged across the grid
(Fig.~\ref{fig:dose-panels}). These patterns resemble the growth of
ordered regions in other interacting systems away from
equilibrium~\cite{hinrichsen2000nonequilibrium,bray1994theory}, and the
dependence on the fraction of traders given the new rule connects the
model to accounts of how local changes spread through a
population~\cite{granovetter1978threshold,watts2002simple,aizenman1988metastability,holroyd2003sharp}.
Larger grids sharpened the rise in majority probability, but the
switched fraction at which half the markets reached a majority changed
with observation time
(Fig.~\ref{fig:finite-size}; Supplementary
Table~\ref{tab:revision-horizons}), so a market that remained mixed at
one endpoint could still reach a majority later. Neither the displayed
markets nor the sampled majority-probability curves established a
threshold for indefinitely large or long-lived
markets~\cite{binder1981finite,landau2021guide}. The resemblance to
droplets forming and joining also did not establish classical
nucleation, because we measured no energy barrier, surface tension,
formation rate, or critical group size separating growth from decay.
Thus, the supported mechanism was the spread of selling through
neighboring decisions on finite grids over finite observation periods.

How far the price fell depended on how the dealer's resources entered
the price. Dealer models commonly connect trading to intermediary cash
and
inventory~\cite{garman1976market,stoll1978supply,ho1981optimal,braunmunzinger2018agent},
but holdings alone do not specify how those resources translate into
prices. Here, the constant-product reserve curve set both the uniform
price at which a batch of orders executed and the quoted price left
after settlement~\cite{angeris2022constant}. When the dealer bought
shares, its share reserve rose and its cash reserve fell, so continued net
selling lowered execution prices. Shallow positive reserves therefore
allowed large proportional declines while rejecting few eligible orders
(Fig.~\ref{fig:reservoir}C,D). Zero dealer depth was a separate case, in
which the dealer took no net position, only opposing trader orders
matched, and the price stayed fixed, so it could not be read as evidence that every small dealer protects the market from a crash. The no-switch
controls mattered as well, because markets with shallow reserves had
large drawdowns even without the rule switch, and only the paired
increases separated that baseline decline from the effect of the switch
(Supplementary Table~\ref{tab:paired-switch}). Thus, the size of a crash
depended on the execution and price mechanism as well as on the
resources available, and scarce resources did not by themselves limit
the decline.

Financial conditions could still change later decisions, but only
through the technical traders, who responded to the quoted price.
Changing dealer depth or trader holdings altered completed orders and
prices, which changed the technical traders' decisions and therefore
the neighboring intentions seen by adjacent interior traders
(Figs.~\ref{fig:reservoir} and~\ref{fig:balance-sheet}). With the
technical traders held, the same financial changes left every interior
decision unchanged, which made the held condition an exact control for
this effect. With technical traders active, the majority probability
neither rose nor fell consistently with dealer depth, even though
drawdown at switched fraction $f=0.20$ fell as positive depth increased.
Technical traders' Buy and Sell thresholds lay around a moving average
price, at distances that normally adapted over time. Fixing those
distances at each market's pre-switch values preserved price feedback,
but the paired comparison with adapting distances was too uncertain to
establish whether adaptation explained majority-probability differences
across dealer depths
(Supplementary Table~\ref{tab:revision-band}). Varying dealer depth changed
funding and initial price sensitivity together, so the depth comparisons
could not separate the two effects (Supplementary
Table~\ref{tab:revision-impact}). Drawdown also depended on our choice of a reserve mechanism; other
models link order flow to prices through dealer inference or order-book
dynamics~\cite{kyle1985continuous,farmer2005predictive,bouchaud2009how}.
Majority probability measured the share of markets reaching a Sell
majority, whereas drawdown measured the largest price decline. Neither
outcome alone determined the other.

The comparison with established agent-based markets concerns how traders
interact and how their rules change: we assigned a persistent rule
switch instead of modeling learning or strategic choice. In the minority
game, traders choose among strategies scored against shared past outcomes
and interact through the market outcome~\cite{challet1997minority,challet2000modeling}.
In the Santa Fe artificial stock market, traders learn forecasting rules
from a shared price and dividend~\cite{lebaron1999time}. Related studies of the Santa Fe market examined incentives to adopt technical
rules~\cite{joshi1999technical} and choose how often to revise
forecasts~\cite{joshi2002financial}.
Models with chartists and fundamentalists let traders move between the
two groups at rates set by recent profits or by others'
opinions~\cite{lux1999scaling,brock1998heterogeneous,kirman1993ants}.
Lattice and cellular-automaton markets encoded local interactions
through stochastic imitation or spin-like responses, sometimes alongside
market-wide
signals~\cite{bornholdt2001ising,iori2002microsimulation,bartolozzi2004stochastic,qiu2007cellular,sornette2006importance,fan2009effect,atman2012influence}.
Wolfram's one-dimensional market and a later model built on Rule~110
used deterministic rules and set the price from the numbers of buyers
and sellers~\cite{wolfram2002new,valsan2022on}. Here rules did not evolve
through learning: the intervention changed one response for a chosen
fraction of traders, and the paired control retained the baseline rule.
Only technical traders along the edge responded to price, allowing the
rule switch to be tested with price feedback active or removed. Orders
needed cash or shares and were completed by other traders or by a dealer
with limited holdings, so a Sell majority and a price crash were measured
separately. These choices gave up the learning and strategic behavior of
established models in exchange for a rule space that can be enumerated
and an intervention that can be repeated exactly. The questions answered
were correspondingly narrower: not whether technical trading emerges or
whether the market becomes efficient, but how many traders had to change
a rule before a Sell majority formed, whether its formation required
feedback from price, and how resources and execution determined the
resulting price decline.

Several modeling choices limit how far these conclusions extend. Traders
followed assigned rules rather than choosing strategies to maximize
returns, and they updated together on a regular grid under one baseline
lookup table in which different neighbor configurations with the same
sum produced the same response. A trader that had sold all its shares
could still influence its neighbors and count toward a Sell majority,
because we assumed that intentions remained observable and influential even when a trader could not act on them. We did not test how the results would
change if a trader's resources also limited the intentions its neighbors
could see. The technical traders followed one band rule and occupied
selected positions along the edge of the grid, so other placements and
price-based strategies remain untested. Unit orders cleared at a common
base price against a constant-product dealer, a specific reserve
mechanism rather than a calibrated representation of an equity
exchange. The dealer collected prescribed fees in the fee comparisons
but did not set its charges or its participation to maximize profit, and
we did not model incentives to provide or withdraw liquidity. The model
had no order book, dividends, borrowing, or strategic risk management by
the dealer, and its steps and price response were not fitted to market
data, so its simulated crashes did not predict the timing or magnitude of actual market crashes. Settlement conserved cash and shares, and an isolated round trip
restoring dealer shares also restored dealer trading cash. Neither
property established economic realism or ruled out profits from other
traders. The long no-switch controls
extended observation of existing starting markets without establishing
their behavior indefinitely (Appendix~\ref{app:revision-controls}), and
neither the baseline period nor an average over final states showed that
the distribution of market states had stopped changing. Other decision
rules, asynchronous updates, alternative interaction networks, and
different execution mechanisms would test how the connection between
local decisions and price declines depends on these assumptions.
